% ESM D -- Computation notes (revision v2: original Appendix A.3 and A.4 plus the H05 gradient-norm ratio definition)
\setcounter{table}{0}\setcounter{figure}{0}
\section{Method details and computation notes}\label{sec:ek-hesap}

This section gives the details moved out of the Methods section of the main article. These are the hypothesis card and decision rules, the expansion of the loss terms, the training recipe and the intervention details. It also gives computation notes. The screening probabilities, the gradient-norm ratio and the sources of the plausibility anchor are in Appendix~\ref{sec:deney-kaydi} of the main article.

\subsection{Hypothesis card and decision rule details}\label{sec:ek-kart}
In addition to the fields described in the main article, the card also fixes the post-acceptance ablations. The card file must be timestamped before the first run. In practice, the evidential strength of this rule is limited (Section~\ref{sec:sinirliliklar} of the main article). A combined card can be written only if each of its components has been accepted separately. For cards for which no isolated VOC ablation could be found, the expected ranges are narrow and centered on zero.

A false elimination at screening is costly. The threshold was therefore kept below the single-seed variability. Even at zero effect, the pass probability is not negligible, so a pass requires a second confirmation. Some cards eliminated at screening were later completed to three seeds. For such a card, failing the KEEP criterion is described as a ``hit'', and meeting it as a ``missed gain''. A hit is a specificity measure. It is likely even under an uninformative gate, so it is not considered informative (Appendix~\ref{sec:ek-onele} of the main article). In the KEEP rule, $s_\Delta$ is the sample s.d. of the $\Delta_{k,s}$. The alternative of the one-sided test is $H_1:\bar\Delta_k>0$. In the INSPECT case, an additional seed is added only from the K5 budget, for both the variant and the baseline. The KEEP rule was frozen in round 2, after the H01--H03 results had been seen. A card that passes the round-2 rule ($\bar\Delta_k\ge+0.40$) also passes the round-1 rule, but not vice versa.

\subsection{Expansion of the baseline loss terms}\label{sec:ek-kayip}
The neck of the baseline is a path aggregation network (PAN). It produces outputs at three scales (strides 8, 16 and 32); see the third-party overview \cite{khanam_2024_yolov11overview}. In TAL, the class target of the positives is $t_i=a_{ij}\max_{i'}u_{i'j}/\max_{i'}a_{i'j}$. Eq.~\eqref{eq:box} of the main article uses CIoU \cite{zheng_2020_diou} and the distribution focal loss (DFL) \cite{li_2020_gfl}. They are defined as follows:
\begin{align}
\mathrm{CIoU}&=\mathrm{IoU}-\frac{\rho^2(\hat{\mathbf m},\mathbf m)}{d^2}-\alpha_v v,\label{eq:ciou}\\
\mathcal L_{\mathrm{dfl}}&=\frac1T\sum_{i\in\mathcal P}w_i\,\frac14\sum_{e=1}^{4}\Bigl[-\bigl((r_e+1)-y_e\bigr)\log S_{e,r_e}-(y_e-r_e)\log S_{e,r_e+1}\Bigr].\label{eq:dfl}
\end{align}
In Eq.~\eqref{eq:ciou}, $\hat{\mathbf m}$ and $\mathbf m$ are the box centers and $\rho$ is the Euclidean distance. $d$ is the diagonal of the smallest box enclosing both boxes. The aspect-ratio term is $v=\frac4{\pi^2}\bigl(\arctan\frac{w^{gt}}{h^{gt}}-\arctan\frac{\hat w}{\hat h}\bigr)^{2}$, and $\alpha_v=v/\bigl((1-\mathrm{IoU})+v\bigr)$ \cite{zheng_2020_diou}. BCE is $\mathrm{BCE}(z,t)=-t\log\sigma(z)-(1-t)\log(1-\sigma(z))$. The DFL in Eq.~\eqref{eq:dfl} \cite{li_2020_gfl} predicts each of the four edge distances with a softmax distribution $S_e$ over $R=16$ discrete bins. Here $y_e$ is the target distance and $r_e=\lfloor y_e\rfloor$.

\subsection{Training recipe and data augmentation defaults}\label{sec:ek-recete}
The transforms are applied in the following order. Mosaic, with probability 1.0, combines four images on a $2s\times2s$ canvas ($s=640$) \cite{bochkovskiy_2020_yolov4}. RandomPerspective draws the scale factor $a\sim U[1-\sigma_s,\,1+\sigma_s]$ with $\sigma_s=0.5$ and applies a translation of 0.1. Rotation, shear and perspective are zero. The output is cropped to the central $s\times s$ window. Next come an HSV shift (hue-saturation-value; 0.015, 0.7, 0.4) and a horizontal flip with probability 0.5. MixUp and the built-in copy-paste are disabled. Mosaic is turned off in the last 10 epochs (\texttt{close\_mosaic=10}). The built-in copy-paste requires masks, so it has no effect on the box-labeled VOC \cite{ultralytics_2026_augdocs}. H03 and H08 therefore rely on a new box-based implementation.

\subsection{Intervention details}\label{sec:ek-mudahale}
\paragraph{Scale and box loss.}
H01 aims to enforce scale invariance by showing the same object at different pixel densities. In H04, the RandomPerspective scale range is reduced from $[0.5,1.5]$ to $[0.75,1.25]$; the amplitude is not increased. In H02, the auxiliary box for each positive anchor is $b^{\mathrm{in}}=(x_c,\,y_c,\,rw,\,rh)$. The DFL term, TAL and the gains are unchanged. The implementation was verified with an $r=1$ identity unit test. In H05, the effective gradient norm is measured (Appendix~\ref{sec:ek-gradyan} of the main article).

\paragraph{Box-based copy-paste.}
Copy-paste tries to close the small-object gap. It replicates small objects within the same image without overlap \cite{kisantal_2019_smallaug}. The transform runs after mosaic and before RandomPerspective. Candidates are objects that do not overlap other boxes ($\mathrm{IoA}\le0.1$). From these, $k\sim U\{1,2,3\}$ objects are selected and pasted with 3-pixel alpha blending. The steps are given in Section~\ref{sec:ek-cp}. A maskless patch also carries a piece of background. The risk that the model learns the patch edge is one reason why AP$_{75}$ is monitored. The transform is disabled in the last 10 epochs. H03 uses $p=0.5$ and $\tau_A=48$ pixels, and H08 uses $p=1.0$ and $\tau_A=96$ pixels. The paste thresholds are measured on the mosaic canvas. The AP$_S$/AP$_M$/AP$_L$ boundaries are measured in original pixels, following the COCO definition \cite{lin_2014_coco}. In H08, enlarging the pool also changes the size mix of the pasted objects. Probability and pool were therefore not separated in H08. The dose $D=N_{\text{pasted}}/N_{\mathrm{mosaic}}$ [objects/mosaic] is measured approximately. It is derived from counters that the worker processes write periodically. The H08 card pre-specified a target dose of $D\ge0.80$, with an expected range of 0.8--1.2. If the $D$ measured before the run lies outside the gate $[0.6,\,1.5]$, the run does not start. A dose below the target is not adjusted, to avoid introducing a second variable. The realized dose is reported.

\paragraph{Training-free inference rows (H07).}
The reference pipeline R0 uses a confidence threshold of 0.001 and an NMS IoU threshold of 0.7. It keeps at most 300 detections and uses a single scale of 640 pixels with class-wise classic NMS. Each row differs from R0 in only one component, and rows are not combined. \emph{NMS threshold (N)} sets the IoU threshold to one of $\{0.5,\,0.6,\,0.8\}$. \emph{Soft-NMS (S)} decays the score of an overlapping box per class with the Gaussian rule, using $\sigma\in\{0.3,\,0.5\}$ \cite{bodla_2017_softnms}. Its definition is given in Section~\ref{sec:ek-h07-dogrulama}. \emph{Test resolution (U)} increases the input size to 704, 768 or 832 pixels. \emph{TTA (T)}: T1 is the Ultralytics default \cite{ultralytics_2026_yolov5tta,ultralytics_2026_v84165}. It runs three branches with scales $\{1,\,0.83,\,0.67\}$ and flips the middle branch horizontally. The branch outputs are passed to a single NMS in the original coordinates. In T2, the scale list is $\{1,\,1.2,\,1.4\}$. The three branches raise the total number of input pixels by a factor of $1+0.83^2+0.67^2\approx2.14$. Latency grows by this factor plus the cost of branch merging and NMS. \emph{WBF (W)} passes the T1 or T2 branches through NMS separately and then merges them by fusion \cite{solovyev_2021_wbf}. It uses an IoU threshold of 0.6, a skip threshold of 0.001 and average confidence. The latency ratio is computed relative to R0 measured in the same session. The threshold $\rho\le1.10$ was set before the run. It is the largest latency increase under which a real-time claim can still be made. GAIN rows do not enter the parent chain of the training cards.


\subsection{Statistical details}\label{sec:ek-istatistik}
This section expands on the ``Statistics'' paragraph in Section~\ref{sec:karar} of the main article. For training cards, seed pairing is only nominal. The intervention also changes how the data augmentation consumes the random number generator (RNG), so pairing reduces the variance only partially. The Holm correction \cite{holm_1979_sequential} uses a K5 family of twelve tests: H03 and TTA, each $\times$ six metrics. This count includes H03 mAP50-95 and is conservative. Holm-7 is applied to the three-seed validation differences of the seven cards. The image-level paired bootstrap \cite{efron_1979_bootstrap} applies the same image indices to the variant and the baseline. It uses $B=1000$ and a 95\% percentile interval on the mean $\Delta$ over the three seeds. This interval does not include training-seed variance. For training interventions, we therefore also give a seed-level interval $\bar\Delta_k\pm t_{0.975,2}\,s_\Delta/\sqrt n$, with $t_{0.975,2}=4.303$. Uncorrected $p<0.05$ is not color-coded in the tables. The practical margins ($+0.30$ and $+0.40$ points) are a criterion separate from the statistical one. In K5, the test mAP50-95 of the baseline is compared, for each seed, with the value from the plausibility check. Any discrepancy is reported without correction.

\subsection{External reference value}\label{sec:ek-dis}
Only one external reference value is available for the order of magnitude of the KITTI baseline. It is a community model card (YOLO11n, KITTI; \url{https://huggingface.co/dronefreak/kitti-yolo11n}, accessed 2026-10-22). The card reports 24.0 for validation mAP50-95 and 67.3 for car AP50-95. The model was trained at 1280 px on the 3712/3769 split. The handling of \texttt{DontCare} is not specified. Early stopping (epoch 43) and the reported numbers rely on the same validation set. Resolution, split and evaluation differ, so these values cannot be compared with ours. The model card is not a peer-reviewed source.

\FloatBarrier
